% GENERATED FILE. Edit canonical lessons, then run npm run generate:books.
% canonical-source-hash: fnv1a64:46b0f2b0211c107b
% canonical-lessons: KA-C536-nataka, KA-C536-rangabhumi, KA-C536-mukhavada, KA-C536-pradarsisu, KA-C536-abhinayisu

\chapter{On Stage}
\label{ch:ka-on-stage}
\hlchaptermodality{536}

\begin{chapteropening}
\textbf{By the end of this chapter:} \emph{I can say a stage play, theatre, a mask, to perform and to act.}

Complete the last lesson of chapter 536: I can say a stage play, theatre, a mask, to perform and to act.
\end{chapteropening}

\section[\texorpdfstring{\kn{ನಾಟಕ} (nāṭaka)}{ನಾಟಕ (nāṭaka)}]{\kn{ನಾಟಕ} (nāṭaka) — a stage play; a drama}

\label{lesson:KA-C536-nataka}



\begin{quote}
Before the new one: say the Kannada for a violin, then the Kannada for to beat a drum.
\end{quote}

\subsection*{You'll want to know: \kn{ನಾಟಕ}}
\textbf{\kn{ನಾಟಕ}} — \emph{nāṭaka} — \textquotedblleft{}a stage play; a drama\textquotedblright{}.

\begin{quote}
\kn{ಶಾಲೆಯ} \kn{ನಾಟಕದಲ್ಲಿ} \kn{ನಾನು} \kn{ರಾಜನ} \kn{ಪಾತ್ರ} \kn{ಮಾಡಿದೆ}.
\end{quote}

\begin{quote}
\emph{śāleya nāṭakadalli nānu rājana pātra māḍide.}
\end{quote}

\begin{quote}
I played the king's part in the school play.
\end{quote}

\kn{ನಾಟಕ} \kn{ಆಡು}, to stage a play.

\begin{cousinweb}
\textbf{\kn{ನಾಟಕ}} (\emph{nāṭaka}) is borrowed from Sanskrit \emph{nāṭaka}, \textquotedblleft{}a dancer, an actor\textquotedblright{}, and later a drama (CDIAL 7045), built on \emph{naṭa}, \textquotedblleft{}an actor\textquotedblright{}.
\end{cousinweb}

\begin{grammarlens}[title={What you've built}]
One more word for describing what you lived and felt.
\end{grammarlens}

\subsection*{Your turn}
\begin{itemize}
\raggedright
  \item \emph{Say it:} \emph{nāṭaka}
  \item \emph{Say it:} \emph{nāṭaka}, once more
  \item \emph{Say it:} \emph{śāleya nāṭakadalli nānu rājana pātra māḍide.} — I played the king's part in the school play.
  \item \emph{Recall:} say the Kannada for a violin, then the Kannada for to beat a drum, then say \emph{nāṭaka} again
\end{itemize}

\begin{tcolorbox}[breakable,colback=teal!4,colframe=teal!35!black,title={Before you move on}]
What does \kn{ನಾಟಕ} mean? (\textquotedblleft{}A stage play; a drama\textquotedblright{}.) Say it once more.
\end{tcolorbox}
\section[\texorpdfstring{\kn{ರಂಗಭೂಮಿ} (raṅgabhūmi)}{ರಂಗಭೂಮಿ (raṅgabhūmi)}]{\kn{ರಂಗಭೂಮಿ} (raṅgabhūmi) — theatre; the stage}

\label{lesson:KA-C536-rangabhumi}



\begin{quote}
Before the new one: say the Kannada for to beat a drum, then the Kannada for a stage play.
\end{quote}

\subsection*{You'll want to know: \kn{ರಂಗಭೂಮಿ}}
\textbf{\kn{ರಂಗಭೂಮಿ}} — \emph{raṅgabhūmi} — \textquotedblleft{}theatre; the stage\textquotedblright{}.

\begin{quote}
\kn{ನನ್ನ} \kn{ಮಾವ} \kn{ರಂಗಭೂಮಿಯ} \kn{ಹಳೆಯ} \kn{ನಟ}.
\end{quote}

\begin{quote}
\emph{nanna māva raṅgabhūmiya haḷeya naṭa.}
\end{quote}

\begin{quote}
My uncle is a veteran stage actor.
\end{quote}

\kn{ರಂಗ}, stage + \kn{ಭೂಮಿ}, ground.

\begin{cousinweb}
\textbf{\kn{ರಂಗಭೂಮಿ}} (\emph{raṅgabhūmi}) is a learned compound built from Sanskrit \emph{raṅga}, \textquotedblleft{}a place of public amusement, a stage\textquotedblright{} (CDIAL 10561), and \emph{bhūmi}, \textquotedblleft{}earth, ground\textquotedblright{} (CDIAL 9557).
\end{cousinweb}

\begin{grammarlens}[title={What you've built}]
One more word for describing what you lived and felt.
\end{grammarlens}

\subsection*{Your turn}
\begin{itemize}
\raggedright
  \item \emph{Say it:} \emph{raṅgabhūmi}
  \item \emph{Say it:} \emph{raṅgabhūmi}, once more
  \item \emph{Say it:} \emph{nanna māva raṅgabhūmiya haḷeya naṭa.} — My uncle is a veteran stage actor.
  \item \emph{Recall:} say the Kannada for to beat a drum, then the Kannada for a stage play, then say \emph{raṅgabhūmi} again
\end{itemize}

\begin{tcolorbox}[breakable,colback=teal!4,colframe=teal!35!black,title={Before you move on}]
What does \kn{ರಂಗಭೂಮಿ} mean? (\textquotedblleft{}Theatre; the stage\textquotedblright{}.) Say it once more.
\end{tcolorbox}
\section[\texorpdfstring{\kn{ಮುಖವಾಡ} (mukhavāḍa)}{ಮುಖವಾಡ (mukhavāḍa)}]{\kn{ಮುಖವಾಡ} (mukhavāḍa) — a mask}

\label{lesson:KA-C536-mukhavada}



\begin{quote}
Before the new one: say the Kannada for a stage play, then the Kannada for theatre.
\end{quote}

\subsection*{You'll want to know: \kn{ಮುಖವಾಡ}}
\textbf{\kn{ಮುಖವಾಡ}} — \emph{mukhavāḍa} — \textquotedblleft{}a mask\textquotedblright{}.

\begin{quote}
\kn{ನಾಟಕದಲ್ಲಿ} \kn{ನಟರು} \kn{ಮುಖವಾಡ} \kn{ಹಾಕಿದ್ದರು}.
\end{quote}

\begin{quote}
\emph{nāṭakadalli naṭaru mukhavāḍa hākiddaru.}
\end{quote}

\begin{quote}
The actors wore masks in the play.
\end{quote}

Also a false face: \kn{ಅವನ} \kn{ಮುಖವಾಡ} \kn{ಕಳಚಿತು}, his mask slipped.

\begin{cousinweb}
\textbf{\kn{ಮುಖವಾಡ}} (\emph{mukhavāḍa}) is built on Sanskrit \emph{mukha}, \textquotedblleft{}mouth, face\textquotedblright{} (CDIAL 10158); the origin of the ending -\emph{vāḍa} is not settled.
\end{cousinweb}

\begin{grammarlens}[title={What you've built}]
One more word for describing what you lived and felt.
\end{grammarlens}

\subsection*{Your turn}
\begin{itemize}
\raggedright
  \item \emph{Say it:} \emph{mukhavāḍa}
  \item \emph{Say it:} \emph{mukhavāḍa}, once more
  \item \emph{Say it:} \emph{nāṭakadalli naṭaru mukhavāḍa hākiddaru.} — The actors wore masks in the play.
  \item \emph{Recall:} say the Kannada for a stage play, then the Kannada for theatre, then say \emph{mukhavāḍa} again
\end{itemize}

\begin{tcolorbox}[breakable,colback=teal!4,colframe=teal!35!black,title={Before you move on}]
What does \kn{ಮುಖವಾಡ} mean? (\textquotedblleft{}A mask\textquotedblright{}.) Say it once more.
\end{tcolorbox}
\section[\texorpdfstring{\kn{ಪ್ರದರ್ಶಿಸು} (pradarśisu)}{ಪ್ರದರ್ಶಿಸು (pradarśisu)}]{\kn{ಪ್ರದರ್ಶಿಸು} (pradarśisu) — to perform; to put on show}

\label{lesson:KA-C536-pradarsisu}



\begin{quote}
Before the new one: say the Kannada for theatre, then the Kannada for a mask.
\end{quote}

\subsection*{You'll want to know: \kn{ಪ್ರದರ್ಶಿಸು}}
\textbf{\kn{ಪ್ರದರ್ಶಿಸು}} — \emph{pradarśisu} — \textquotedblleft{}to perform; to put on show\textquotedblright{}.

\begin{quote}
\kn{ಮಕ್ಕಳು} \kn{ವೇದಿಕೆಯ} \kn{ಮೇಲೆ} \kn{ಒಂದು} \kn{ಸಣ್ಣ} \kn{ನಾಟಕ} \kn{ಪ್ರದರ್ಶಿಸಿದರು}.
\end{quote}

\begin{quote}
\emph{makkaḷu vēdikeya mēle ondu saṇṇa nāṭaka pradarśisidaru.}
\end{quote}

\begin{quote}
The children performed a short play on the stage.
\end{quote}

Past (I …): \kn{ಪ್ರದರ್ಶಿಸಿದೆ}, \emph{pradarśiside}.

Also to exhibit: \kn{ಚಿತ್ರಗಳನ್ನು} \kn{ಪ್ರದರ್ಶಿಸು}. \kn{ಪ್ರದರ್ಶನ} is a show.

\begin{cousinweb}
\textbf{\kn{ಪ್ರದರ್ಶಿಸು}} (\emph{pradarśisu}) is built on Sanskrit \emph{pradarśana}, \textquotedblleft{}a showing, look, appearance\textquotedblright{} (Monier-Williams), with the verb ending -\emph{isu}.
\end{cousinweb}

\begin{grammarlens}[title={What you've built}]
One more word for describing what you lived and felt.
\end{grammarlens}

\subsection*{Your turn}
\begin{itemize}
\raggedright
  \item \emph{Say it:} \emph{pradarśisu}
  \item \emph{Say it:} \emph{pradarśisu}, once more
  \item \emph{Say it:} \emph{makkaḷu vēdikeya mēle ondu saṇṇa nāṭaka pradarśisidaru.} — The children performed a short play on the stage.
  \item \emph{Recall:} say the Kannada for theatre, then the Kannada for a mask, then say \emph{pradarśisu} again
\end{itemize}

\begin{tcolorbox}[breakable,colback=teal!4,colframe=teal!35!black,title={Before you move on}]
What does \kn{ಪ್ರದರ್ಶಿಸು} mean? (\textquotedblleft{}To perform; to put on show\textquotedblright{}.) Say it once more.
\end{tcolorbox}
\section[\texorpdfstring{\kn{ಅಭಿನಯಿಸು} (abhinayisu)}{ಅಭಿನಯಿಸು (abhinayisu)}]{\kn{ಅಭಿನಯಿಸು} (abhinayisu) — to act}

\label{lesson:KA-C536-abhinayisu}



\begin{quote}
Before the new one: say the Kannada for a mask, then the Kannada for to perform.
\end{quote}

\subsection*{You'll want to know: \kn{ಅಭಿನಯಿಸು}}
\textbf{\kn{ಅಭಿನಯಿಸು}} — \emph{abhinayisu} — \textquotedblleft{}to act\textquotedblright{}.

\begin{quote}
\kn{ಆ} \kn{ಸಿನಿಮಾದಲ್ಲಿ} \kn{ನನ್ನ} \kn{ಸ್ನೇಹಿತ} \kn{ಅಭಿನಯಿಸಿದನು}.
\end{quote}

\begin{quote}
\emph{ā sinimādalli nanna snēhita abhinayisidanu.}
\end{quote}

\begin{quote}
My friend acted in that film.
\end{quote}

Past (I …): \kn{ಅಭಿನಯಿಸಿದೆ}, \emph{abhinayiside}.

\kn{ಅಭಿನಯ} is acting.

\begin{cousinweb}
\textbf{\kn{ಅಭಿನಯಿಸು}} (\emph{abhinayisu}) is built on Sanskrit \emph{abhinaya}, from \emph{abhi-nī}, \textquotedblleft{}to bring near, to act out\textquotedblright{} (Monier-Williams), with the verb ending -\emph{isu}.
\end{cousinweb}

\begin{grammarlens}[title={What you've built}]
One more word for describing what you lived and felt.
\end{grammarlens}

\subsection*{Your turn}
\begin{itemize}
\raggedright
  \item \emph{Say it:} \emph{abhinayisu}
  \item \emph{Say it:} \emph{abhinayisu}, once more
  \item \emph{Say it:} \emph{ā sinimādalli nanna snēhita abhinayisidanu.} — My friend acted in that film.
  \item \emph{Recall:} say the Kannada for a mask, then the Kannada for to perform, then say \emph{abhinayisu} again
\end{itemize}

\begin{tcolorbox}[breakable,colback=teal!4,colframe=teal!35!black,title={Before you move on}]
What does \kn{ಅಭಿನಯಿಸು} mean? (\textquotedblleft{}To act\textquotedblright{}.) Say it once more.
\end{tcolorbox}
